% !TEX root = sga4.5-en.tex
% ENGLISH EDITION -- TRANSLATED FROM FRENCH (AI-assisted, needs human review).


\chapter{Finiteness theorems in \texorpdfstring{$\ell$}{l}-adic cohomology}\label{VII}











\section{Statement of the theorems}\label{VII:1}

Throughout this expos\'e, $S$ will be a noetherian scheme and $A$ a
left-noetherian ring, with torsion annihilated by an integer invertible on $S$.
Our main result is the following.




\begin{theorem_}\label{VII:1-1}
Assume $S$ regular of dimension $0$ or $1$. Let $f:X\to Y$ be a
morphism of $S$-schemes of finite type and $\sF$ a constructible sheaf of
left $A$-modules on $X$. Then the sheaves $\eR^i f_\ast \sF$ are
constructible.
\end{theorem_}

The proof will be given in paragraph \ref{VII:2} and in \ref{VII:3-10}.




\subsection{Remark}\label{VII:1-2}

In characteristic $0$, this result is less general than
\cite[XIX paragraph 5]{sga4}, which proves the conclusion of the theorem for
every morphism of finite type of excellent schemes of characteristic $0$.
The proof of \cite[XIX]{sga4} uses on the one hand resolution of
singularities, on the other hand that the schemes be of equal
characteristic (in order to deduce from the resolution the ``purity
theorem'').




\subsection{Remark}\label{VII:1-3}

We say that a complex $K\in\ob\eD(X,A)$ is \emph{constructible} if its
cohomology sheaves are constructible and we denote with a subscript
$c$ the subcategory of $\eD(X,A)$ (or $\eD^+$, $\eD^-$, $\eD^b$) formed
of constructible complexes. In the language of derived categories,
which we shall use freely, \ref{VII:1-1} says that
$\eR f_\ast:\eD^+(X,A) \to \eD^+(Y,A)$ sends $\eD_c^+$ into $\eD_c^+$.
Explicitly:
\begin{enumerate}[\indent a)]
  \item For $K$ reduced to a sheaf $\sF$ in degree $0$, one has
    $\sH^i\eR f_\ast K = \eR^i f_\ast \sF$; the derived statement therefore
    implies the theorem.
  \item In the other direction, one invokes the spectral sequence
    \[
      E_1^{p q} = \eR^q f_\ast \sH^p(K) \Rightarrow \sH^{p+q} \eR f_\ast K \text{.}
    \]
\end{enumerate}

Speaking the language of derived categories has the advantage of replacing by
a simple transitivity formula $\eR(f g)_\ast = \eR f_\ast \eR g_\ast$
a Leray spectral sequence
$E_2^{p q} = \eR^p f_\ast \eR^q g_\ast \Rightarrow \eR^{p+q} (f g)_\ast$.

Under the hypotheses of \ref{VII:1-1}, it follows from \cite{sga4} that
$\eR f_\ast$ is of finite cohomological dimension. This allows above to
replace $\eD^+$ by $\eD$, $\eD^-$ or $\eD^b$.




\subsection{The ideas of the proof}\label{VII:1-4}
\begin{enumerate}[\indent a)]
  \item Poincar\'e duality allows one to treat the case where $X$ is
    smooth, $\sF$ locally constant, and $Y=S$. The dimension hypothesis
    appears when computing $\rHom$'s over $S$.
  \item We factor $f$ as $g j$: $g$ proper and $j$ an open immersion. One has
    $\eR f_\ast = \eR g_\ast \eR j_\ast$ and the finiteness theorem for
    proper morphisms controls $\eR g_\ast$. By d\'evissage, we reduce to
    supposing $x$ smooth, $\sF$ locally constant, $f$ an open immersion
    and $Y$ proper over $S$.
  \item Induction on $\dim X$ (with change of $S$) allows, roughly
    speaking, to suppose the theorem true outside a part of $Y$ finite
    over $S$. Denoting by $a:X\to S$ and $b:Y\to S$ the structural morphisms, one
    then shows that if the theorem were false for $\sF$ and
    $\eR f_\ast$, it would also be false for $\sF$ and
    $\eR b_\ast \eR f_\ast = \eR a_\ast$: contradiction.
\end{enumerate}




\begin{corollary_}\label{VII:1-5}
Under the hypotheses of the theorem, the categories $\eD_c(X,A)$ and
$\eD_c(Y,A)$ are interchanged by the $4$ operations
$\eR f_\ast$, $\eR f_!$, $f^\ast$, $\eR f^!$.
\end{corollary_}

$\eR f_\ast$ is treated above, $\eR f_!$ in \cite[XVII 5.3]{sga4},
$f^\ast$ is clear. There remains $\eR f^!$. The problem is local, which allows us to
assume that $f$ factors as $X\xrightarrow i Z \xrightarrow g Y$ ($i$
a closed immersion and $g$ smooth, purely of relative dimension $n$). The
Poincar\'e duality $\eR g_! K(n)[2 n]$ and transitivity
$\eR f^! = \eR i^! \eR g^!$ reduce us to proving \ref{VII:1-5} for $i$.
For every $L\in \eD(Z,A)$, if $j$ is the inclusion in $Z$ of $U=Z\setminus X$,
$i_\ast \eR i^! L$ is the mapping cylinder of $K\to \eR j_\ast j^\ast K$ and
\ref{VII:1-5} for $i$ follows from \ref{VII:1-1} for $j$.




\begin{corollary_}\label{VII:1-6}
Let $X$ be as in the theorem, and assume $A$ commutative. Then, if
$\sF$ and $\sG$ are constructible sheaves of $A$-modules on $X$, the
$\Ext_A^i(\sF,\sG)$ are constructible: $\rHom$ sends
$\eD_c^-(X,A)\times \eD_c^+(X,A)$ into $\eD_c^+(X,A)$.
\end{corollary_}

By d\'evissage of $\sF$, we reduce to assuming $\sF$ of the form
$j_!\sF_1$ ($j:Y\hookrightarrow X$) a locally closed immersion and $\sF_1$
locally constant on $Y$). One then has
\[
  \rHom(j_! \sF_1,\sG) = \eR j_\ast \Hom(\sF_1,\eR j^! \sG) \text{:}
\]
by \ref{VII:1-3}, we are reduced to the case where $\sF$ is locally
constant --- indeed constant since the problem is local. For $\sF$
locally constant constructible, one has $\Hom(\sF,\sG)_x=\hom(\sF_x,\sG_x)$,
and likewise for $\rHom$. For $\sF$ constant, one can therefore compute the
$\Ext^i$ using a finite-type projective resolution of its constant
value, and the $\Ext^i$ are constructible if $\sG$ is --- whence the
corollary.




\subsection{Remark}\label{VII:1-7}

Since $\eR f_\ast$ is of finite cohomological dimension, it transforms
complexes of finite tor-dimension (resp.\ $\leqslant d$) into complexes of
finite tor-dimension (resp.\ $\leqslant d$) \cite[XVII 5.2.11]{sga4}.
\cite[XVII 5.2.10]{sga4} and the proof of \ref{VII:1-5} then show that
finite tor-dimension is stable under the $4$ operations $\eR f_\ast$,
$\eR f_!$, $f^\ast$, $\eR f^!$. A variant of that of \ref{VII:1-6}
(unscrewing $K$ according to a partition of $X$ to assume the cohomology
sheaves locally constant, then localizing to replace $K$ by a
finite complex of locally free sheaves of finite type) then shows that,
for $A$ commutative, $\rHom$ induces
\[
  \rHom:\eD_\text{ctf}^b(X,A)\times \eD_\text{tf}^b(X,A) \to \eD_\text{tf}^+(X,A)
\]
(t.f.: finite tor-dimension).




\subsection{}\label{VII:1-8}

Under the hypotheses of the theorem, the same method allows us to
prove a local biduality theorem $K\iso D D K$ (paragraph
\ref{VII:4}). We also prove a finiteness theorem for the
sheaves of vanishing cycles. For arbitrary $S$, one still obtains a
``generic'' theorem:




\begin{theorem_}\label{VII:1-9}
Let $f:X\to Y$ be a morphism of $S$-schemes of finite type and $\sF$ a
constructible sheaf of $A$-modules on $X$. There exists a dense open $U$ of $S$ such
that
\begin{enumerate}[\indent (i)]
  \item Above $U$, the $\eR^i f_\ast \sF$ are constructible, and zero
    except for finitely many of them.
  \item The formation of the $\eR^i f_\ast \sF$ is compatible with every
    base change $S' \to U\subset S$.
\end{enumerate}
\end{theorem_}

For $S$ the spectrum of a field, one has $U=S$.

If $S$ is the spectrum of a field, a passage-to-the-limit argument provides
compatibility of the $\eR^i f_\ast$ with $S$-base changes for every
quasi-compact quasi-separated morphism of $S$-schemes and every sheaf
$\sF$.




\begin{corollary_}\label{VII:1-10}
Let $x$ be a scheme of finite type over $k$ separably closed and $\sF$ a
constructible sheaf of $A$-modules. Then the $\h^i(X,\sF)$ are of finite
type.
\end{corollary_}

This is the special case $S=\spec(k)$, $Y=S$.

\begin{corollary_}\label{VII:1-11}
Let $X$ and $Y$ be two schemes of finite type over $k$ separably closed and
$K\in \ob\eD^-(X,A^\circ)$, $L\in \ob\eD^-(Y,A)$ complexes of sheaves of
right and left modules. The K\"unneth arrow
\[
  \eR \Gamma(X,K) \lotimes \eR \Gamma(Y,L) \to \eR\Gamma(X\times Y,\pr_1^\ast K \lotimes \pr_2^\ast L)
\]
is an isomorphism.
\end{corollary_}

One proceeds as in \cite[XVII 5.4.3]{sga4}:
\[\xymatrix{
  X\times Y \ar[r]^-{\text{pr}_1} \ar[d]^-{\text{pr}_2}
    & X \ar[d]^-a \\
  Y \ar[r]^-b
    & \spec(k)
}\]
changing base by $b$, one finds that
$\eR\pr_{2,\ast}{} \pr_1^\ast K$ is
$b^\ast \eR\Gamma(X,K)$. One then has (cf.\ \cite[XVII 5.2.11]{sga4} and proof)

\begin{align*}
  \eR\Gamma(X\times Y,\pr_1^\ast K\lotimes \pr_2^\ast L)
    &= \eR\Gamma(Y,\eR \pr_{2,\ast}(\pr_1^\ast K\lotimes \pr_2^\ast L)) \\
    &= \eR\Gamma(Y,(\eR\pr_{2,\ast}\pr_1^\ast K)\lotimes L) \\
    &= \eR\Gamma(Y,b^\ast \eR\Gamma(X,K)\lotimes L) \\
    &= \eR\Gamma(X,K)\lotimes \eR\Gamma(Y,L)
\end{align*}




\subsection{}\label{VII:1-12}

Finally, as a local counterpart of \ref{VII:1-9}(ii), we shall prove a
generic local acyclicity theorem (\ref{VII:2-13}).










\section{Generic theorems}\label{VII:2}

In this paragraph, we prove \ref{VII:1-9} and the theorem
of generic local acyclicity \ref{VII:2-13}. To prove \ref{VII:1-9},
one immediately reduces to assuming $S$ integral. Let $\eta$ be its
generic point.




\subsection{}\label{VII:2-1}

We shall begin by proving \ref{VII:1-9} under the following
additional hypotheses: $X$ is smooth over $S$, purely of relative dimension
$n$, $A=\dZ/m$, $\sF$ is locally constant and $Y=S$. Let
$\sF'=\Hom(\sF,\dZ/m)$. We shall show that if the $\eR^i f_! \sF'$ are
locally constant, the conclusions of \ref{VII:1-9} hold for $U=S$.
Recall that if $\sL$ is a locally constant constructible sheaf, the
$\Ext^i(\sL,\sG)$ are computed fibre by fibre, and are constructible if $\sG$
is. Recall also that $\dZ/m$ is an injective $\dZ/m$-module, and is the
dualizing module, so that $\sF=\rHom(\sF',\dZ/m)$. Poincar\'e duality
\[
  \eR f_\ast \rHom(K,\eR f^! L) = \rHom(\eR f_! K,L) \text{,}
\]
for $K=\sF'$, $L=\dZ/m$, $\eR f^! L = \dZ/m[2n](n)$, thus provides
\[
  \eR^{2n-i} f_\ast \sF = \Hom(\eR^i f_! \sF',\dZ/m)(-n) \text{.}
\]

The sheaf on the right-hand side is locally constant, of formation compatible
with every base change --- whence the assertion.




\subsection{}\label{VII:2-2}

Let us prove \ref{VII:1-9} under the additional hypotheses: $X$ is smooth over
$S$, $\sF$ is locally constant and $Y=S$.
\begin{enumerate}[a)]
  \item Decomposing $X$ into connected components, we reduce to assuming
    that it is purely of some relative dimension $n$ over $S$.
  \item Decomposing $A$ into a product, we reduce to assuming that
    $\ell^m A=0$, with $\ell$ prime invertible on $S$; we then filter $\sF$
    by the $\ell^k \sF$ to reduce, via the corresponding spectral
    sequence, to the case where $\ell\sF=0$. One can then replace $A$ by
    $A/\ell A$ and assume that $\ell A=0$.
  \item Replacing $S$ by a suitable $U$, we may assume that there exists a
    finite \'etale Galois covering $X_1/X$, with Galois group $G$, such
    that the inverse image $\sF_1$ of $\sF$ on $X_1$ is a constant sheaf,
    of constant value $F$. Denoting by $f_1$ the projection of $X_1$ onto
    $S$, one then has
    \[
      \eR^i {f_1}_\ast \sF_1 = (\eR^i {f_1}_\ast \dZ/\ell)\otimes_{\dZ/\ell} F \text{.}
    \]
    One also has the Hochschild-Serre spectral sequence (or: Leray
    for the covering $X_1/X$)
    \[
      \underline\h^p(G,\eR^q {f_1}_\ast \sF_1) \Rightarrow \eR^{p+q} f_\ast \sF \text{.}
    \]
    By \ref{VII:2-1}, we may assume, after shrinking $U$, that
    the $\eR^i {f_1}_\ast \dZ/\ell$ are locally constant, of formation
    compatible with every base change. The same property then holds
    for the $\eR^i f_\ast \sF$. Finally, if $\bar\eta$ is a
    geometric point localized at the generic point $\eta$ of $S$, the
    $\eR^i f_\ast\sF$ are almost all zero, since the
    $(\eR^i f_\ast \sF)_{\bar\eta} = \h^i(X_{\bar \eta},\sF)$ are.
\end{enumerate}




\subsection{}\label{VII:2-3}

Let us prove by induction on $n$ that
\begin{equation*}\tag{$\ast_n$}\label{VII:eq:2-3-1}
  \begin{array}{c}
    \text{The conclusions of \ref{VII:1-9} are true when $\dim{X_\eta}\leqslant n$} \\
    \text{and $f$ is a dense open immersion.}
  \end{array}
\end{equation*}

For $n=0$, after shrinking $S$, one has $X=Y$: $(\ast_0)$ is clear.
Assume ($\ast_{n-1})$, and let us prove \eqref{VII:eq:2-3-1}. In
($\ast_{n-1}$), one may replace ``open immersion'' by ``immersion''
as seen by factoring into an open immersion and a closed immersion.




\begin{lemma_}\label{VII:2-4}
After shrinking $S$, the conclusions of \ref{VII:1-9} hold
above $Y'\subset Y$, the complement $Y_1$ of $Y'$ being finite over $S$.
\end{lemma_}

The assertion is local on $Y$, which we may assume affine: $Y\subset \dA_S^n$.
The induction hypothesis ($\ast_{n-1}$) applies to
\[\xymatrix{
  X \ar[r]^-f \ar[dr]
    & Y \ar[d]^-{\pr_i} \\
  & \dA_S^1
}\]
There therefore exists for each $i$ a dense open $U_i$ of $\dA_S^1$ such that the
conclusions of \ref{VII:1-9} hold above $\pr_i^{-1}(U_i)$; they hold
above the union of the $\pr_i^{-1}(U_i)$, and \ref{VII:2-4} follows.




\addtocounter{subsection}{1} % a section is skipped in the original
\subsection{}\label{VII:2-6}

Let us prove \eqref{VII:eq:2-3-1} for $X$ smooth over $S$ and $\sF$ locally
constant on $X$. The problem being local on $Y$, we may assume $Y$
affine, then projective (replace $Y$ by its closure in a projective
space). Let $i:Y_1\hookrightarrow Y$ and $j:Y' \to Y$ be those guaranteed by
\ref{VII:2-4}.
\[\xymatrix{
  X \ar@{^{(}->}[r] \ar[dr]_-a
    & Y \ar[d]^-b
    & \ar[l]_-i \ar[dl]^-{b_1} Y_1 \\
  & S
}\]
after shrinking $S$, we know that $j^\ast \eR f_\ast \sF$ is
constructible, of formation compatible with every base change in $S$; we
also know that $\eR a_\ast \sF = \eR b_\ast \eR f_\ast \sF$ is constructible,
of formation compatible with every base change.

Applying $\eR b_\ast$ to the triangle defined by the exact sequence
\begin{equation*}\tag{1}\label{VII:eq:2-6-1}
\xymatrix{
  0 \ar[r]
    & j_! j^\ast \eR f_\ast \sF \ar[r]
    & \eR f_\ast \sF \ar[r]
    & i_! i^\ast \eR f_\ast \sF \ar[r]
    & 0 \text{:}
}
\end{equation*}
one obtains a triangle
\begin{equation*}\tag{2}\label{VII:eq:2-6-2}
\xymatrix{
  \ar[r]
    & \eR b_\ast j_! j^\ast \eR f_\ast \sF \ar[r]
    & \eR a_\ast \sF \ar[r]
    & {b_1}_\ast i^\ast \eR f_\ast \sF \ar[r]
    &
}
\end{equation*}
in which the first two terms are constructible, of formation
compatible with every base change in $S$ (for the first, by the
finiteness theorem for the proper morphism $b$). The same therefore holds
for the third. Since $b_1$ is finite, we deduce that
$i^\ast \eR f_\ast \sF$ is constructible of formation compatible with every
base change in $S$, and likewise for $\eR f_\ast \sF$ by
\eqref{VII:eq:2-6-1}.




\subsection{}\label{VII:2-7}

Let us prove \eqref{VII:eq:2-3-1} in general. We begin by reducing to the case
where in $X$ there exists a dense open $V$ smooth over $S$. For $S$ the spectrum of a
perfect field, it suffices to replace $X$ by $X_\text{red}$ (and $Y$ by
$Y_\text{red}$). In general, one must shrink $S$, perform a finite radicial surjective base
change $S' \to S$, and replace $X$ and $Y$ by
$X_\text{red}$ and $Y_\text{red}$. Since the \'etale topology is insensitive to
finite surjective radicial morphisms, this is harmless. After
shrinking $V$, we may assume $\sF$ locally constant on $V$
\[\xymatrix{
  V \ar@{^{(}->}[r]^i
    & X \ar@{^{(}->}[r]^-f
    & Y \text{.}
}\]
Define $\Delta$ by the triangle
\begin{equation*}\tag{1}\label{VII:eq:2-7-1}
\xymatrix{
  \ar[r]
    & \sF \ar[r]
    & \eR j_\ast j^\ast \sF \ar[r]
    & \Delta \ar[r]
    & \text{.}
}
\end{equation*}
The cohomology sheaves of $\Delta$ are supported on $X\setminus V$,
and $\dim(X\setminus V)_\eta < n$. The induction hypothesis therefore allows us to
assume that $\eR f_\ast \Delta$ is constructible. Applying $\eR f_\ast$ to the
triangle \eqref{VII:eq:2-7-1}; one finds a triangle
\[\xymatrix{
  \ar[r]
    & \eR f_\ast \sF \ar[r]
    & \eR(f j)_\ast j^\ast \sF \ar[r]
    & \eR f_\ast \Delta \ar[r]
    &
}\]
in which two of the vertices are constructible of formation compatible with
every base change in $S$. The third is therefore as well.




\subsection{}\label{VII:2-8}

Let us prove \ref{VII:1-9}. The problem is local on $Y$, which we may assume
affine. Taking an affine covering of $X$ and invoking the Leray spectral
sequence, we reduce to having $X$ affine as well. All this to
ensure that we can factor $f$ as an open immersion followed by a
proper morphism: $f=g j$, whence $\eR f_\ast = \eR g_\ast \eR j_\ast$. The
open immersion is amenable to \eqref{VII:eq:2-3-1}, the proper morphism
to the finiteness theorem.

The following corollaries are proved as in paragraph 1.




\begin{corollary_}\label{VII:2-9}
Under the hypotheses of the theorem, for $K$ in $\eD_c^b(X,A)$ or
$\eD_c^b(Y,A)$ respectively, there exists a non-empty open $U$ of $S$
above which $\eR f_\ast K$, $\eR f_! K$, $f^\ast K$, $\eR f^! K$ are
in $\eD_c^b(Y,A)$ or $\eD_c^b(X,A)$, and of formation compatible with every
base change $S' \to U\subset S$.
\end{corollary_}




\begin{corollary_}\label{VII:2-10}
Let $X$ be as in the theorem, and assume $A$ commutative. Then, if
$\sF$ and $\sG$ are constructible sheaves of $A$-modules on $X$,
above a dense open $U$ of $S$, the $\Ext_A^i(\sF,\sG)$ are still
constructible, of formation compatible with every base change
$S' \to U\subset S$.
\end{corollary_}




\subsection{}\label{VII:2-11}

If $x$ is a geometric point of a scheme $X$, we shall denote by $X_x$
the strict henselization of $X$ at $x$. For $f:X\to S$ and $t$ a
geometric point of $S$, we shall denote by $X_t$ the geometric fibre of $X$ at
$t$. Finally, for $x$ a geometric point of $X$, and $t$ a geometric
point of $S_{f(x)}$, $(X_x)_t$ is the fibre at $t$ of
$X_x \to S_{f(x)}$.




\begin{definition_}\label{VII:2-12}
Let $f:X\to S$ and $K\in \ob\eD^+(X,A)$. We say that $f$ is \emph{locally
acyclic at $x$, resp.\ $K$}, if for every geometric point $t$ of
$S_{f(x)}$ one has $K_x=\eR\Gamma(X_x,K)\iso \eR\Gamma(X_{x,t},K)$. We say that $f$
is \emph{locally acyclic, rel.\ $K$}, if this is true for every $x$, and
\emph{universally locally acyclic, rel.\ $K$}, if this remains true
after every base change $S' \to S$.
\end{definition_}




\begin{theorem_}\label{VII:2-13}
Let $f:X\to S$ be a morphism of finite type and $\sF$ constructible on $X$. There
exists a dense open $U$ of $S$ above which $f$ is universally
locally acyclic, rel.\ $\sF$.
\end{theorem_}

We shall admit the following result, proved in the appendix.




\begin{lemma_}\label{VII:2-14}
Let $X\xrightarrow f S_1 \xrightarrow g S_2$. If $g$ is smooth, and $f$
is universally locally acyclic rel.\ $K$, then $g f$ is as well.
\end{lemma_}

We reduce to assuming $S$ integral, with generic point $\eta$, and
proceed by induction on $\dim{X_\eta}$. We begin by deducing from the
induction hypothesis that
\begin{equation*}\tag{A}\label{VII:eq:2-14-1}
  \begin{array}{c}
    \text{After shrinking $S$, there exists $T\subset X$, finite over $S$, such that $f$} \\
    \text{is universally locally acyclic outside $T$.}
  \end{array}
\end{equation*}

This question being local, we localize on $X$ and factor $f$ as
$X\xrightarrow u \dA_S^1 \xrightarrow S$. The induction hypothesis
applies to $u$ and we conclude by \ref{VII:2-14} and the usual arguments.

To prove the theorem, we may assume $X$ proper, since the
problem is local. We shrink $S$ so that the $\eR^i f_\ast \sF$
are locally constant and \eqref{VII:eq:2-14-1} is applicable, and we
use the following lemma.




\begin{lemma_}\label{VII:2-15}
Let $f:X\to S$, proper, and $\sF$ constructible. Assume that the
$\eR^i f_\ast \sF$ are locally constant and that, outside $T\subset X$
finite over $S$, $f$ is locally acyclic rel.\ $\sF$. Then $f$ is
locally acyclic rel.\ $\sF$.
\end{lemma_}

Let $s$ be a geometric point of $S$, $S_s$ the strict localization of $S$
at $s$ and $t$ a geometric point of $S_s$. Let $X_{(s)}$ be the
inverse image of $X$ over $S_s$, $\bar j:X_t \to X_{(s)}$ and
$i:X_s \hookrightarrow X_{(s)}$. Denoting still by $\sF$ the inverse image of
$\sF$ on $X_{(s)}$ or $X_t$, one must prove that
$i^\ast \sF \iso i^\ast \eR \bar j_\ast \sF$. Let $\Delta$ be the mapping cylinder
of this map. Its cohomology sheaves are supported on $T_s$.
To prove that $\Delta=0$, it therefore suffices to prove that
$\eR\Gamma(X_s,\Delta)=0$. This is the mapping cylinder of
$(\eR f_\ast \sF)_s = \eR\Gamma(X_s,\sF) \to \eR\Gamma(X_s,i^\ast \eR\bar j_\ast \sF) = \eR\Gamma(X_{(s)},\eR\bar j_\ast \sF) = \eR\Gamma(X_t,\sF) = (\eR f_\ast \sF)_t$.
This specialization morphism is by hypothesis an isomorphism, and
$\Delta=0$.




\begin{corollary_}\label{VII:2-16}
For $S$ the spectrum of a field, every $S$-scheme $X$ is universally
locally acyclic, rel.\ $K$, whatever $K\in\ob\eD^+(X,A)$.
\end{corollary_}

Follows from \ref{VII:2-13} by passage to the limit (possible, since the $U$ of
\ref{VII:2-13} is always equal to $S$.










\section{Proof of \texorpdfstring{\ref{VII:1-1}}{1.1} and constructibility of sheaves of vanishing cycles}\label{VII:3}




\subsection{}\label{VII:3-1}

Let $S$ be a strictly local trait (the spectrum of a henselian discrete
valuation ring with separably closed residue field). We denote by $s$ and
$\eta$ its closed and generic points, and by $\bar\eta$ a geometric
point localized at $\eta$ (a separable closure of $k(\eta)$).
Let $X$ over $S$ and $\sF$ a sheaf on $X_\eta$. Recall the
definition of the sheaves of vanishing cycles $\eR^i \Psi_\eta(\sF)$ (sheaves
on $X_s$, equipped with an action of $\gal(\bar\eta/\eta)$). Let
$i:X_s\hookrightarrow X$, $j:X_\eta \hookrightarrow X$, and $\bar j$ the
composite $X_{\bar\eta}\to X_\eta \hookrightarrow X$; one has
\[
  \eR^i \Psi_\eta(\sF) = i^\ast \eR^i \bar j_\ast (\bar j^\ast \sF) \text{.}
\]

Here is the main theorem of this number. It improves
\cite[XIII 2.3.1, 2.4.2]{sga7}.




\begin{theorem_}\label{VII:3-2}
Assume $X$ of finite type over $S$ and $\sF$ constructible. Then the
sheaves of vanishing cycles $\eR^i \Psi_\eta(\sF)$ are constructible.
\end{theorem_}

We proceed by induction on the dimension $n$ of $X_\eta$. Moreover we may
suppose --- and we suppose --- $X_\eta$ dense in $X$ (replacing $X$
by $\bar X_\eta$ does not change the $\eR^i \Psi$, which are supported on
$\bar X_\eta$).




\begin{lemma_}\label{VII:3-3}
If \ref{VII:3-2} is true in dimension of $X_\eta<n$, and $\dim X_\eta=n$,
there exist constructible subsheaves $\sG_i$ of the $\eR^i \Psi_\eta(\sF)$
such that the supports of the local sections of $\eR^i\Psi_\eta(\sF)/\sG$ are
finite.
\end{lemma_}

Let $s'$ be a generic geometric point of the affine line over $s$,
and let $S'$ be the strict localization at $s'$ of the affine line $\dA_S^1$ over
$S$. It is again a strictly local trait, and uniformizers for $S$
are uniformizers over $S'$.
\[\xymatrix{
  & \spec(k') \ar[rr] \ar[dd] \ar[dr]
    & & (S',\eta',s') \ar[dr] \ar[dd] \\
  \bar\eta' \ar[ur] \ar[dr]
    & & \dA_{\bar S}^1 \ar[d]
    & & \dA_S^1 \ar[dl] \\
  & \bar\eta\ar@{^{(}->}[r]
    & \bar S \ar[r]
    & (S,\eta,s)
}\]
Let $k'=k(\bar\eta)\otimes_{k(\eta)}k(\eta')$. If $\bar S$ is the normalization
of $S$ in $\bar\eta$, $k'$ is the fraction field of the strict localization
of $\dA_{\bar S}^1$ at $s'$; it is therefore a field, and
$\gal(\bar\eta/\eta)=\gal(k'/\eta')$. Let $\bar\eta'$ be the spectrum of an
algebraic closure of $k'$. The Galois group $P=\gal(\bar\eta'/k')$ is
a pro-$p$-group, for $p$ the characteristic exponent of $k(s)$.




\begin{lemma_}\label{VII:3-4}
For every scheme $X'$ over $S'$, and every sheaf $\sF$ on
$X_\eta'=X_{\eta'}'$, one has between the sheaves of vanishing cycles for
$X'/S'$ and for $X'/S$ the relation
$\eR^i \Psi_\eta(\sF) =\eR^i\Psi_{\eta'}(\sF)^P$.
\end{lemma_}

The following diagram compares the morphisms used in the definition of
$\eR\Psi$ for $X'/S'$ and $X'/S$:
\[\xymatrix{
  X_{\bar\eta'} \ar[rr] \ar[d]
    & & X_{\eta'}' \ar@{=}[dd] \ar@{^{(}->}[dr]
    & & S' \\
  X_{\eta'}\times_{\eta'}\spec(k') \ar@{=}[d] \ar[urr]
    & & & X \ar[ur] \ar[dr] \\
  X_\eta' \ar[rr]
    & & X_\eta' \ar@{^{(}->}[ur]
    & & S
}\]
and the lemma follows from the Hochschild-Serre spectral sequence, taking into
account that $P$ is a pro-$p$-group, with $p$ invertible in $A$.

Let us prove \ref{VII:3-3}. The question is local; this allows us to assume $X$
affine, $X\subset \dA_S^n$. Let $f$ be one of the projections
$X\subset \dA_S^n\to \dA_S^1$. Let $X'$ be the ``localization'' $X_{\dA_S^1} S'$ of $x$,
and $\sF'$ the inverse image of $\sF$ on $x'$
\[\xymatrix{
  \sF'\text{ on }X' \ar[r] \ar[d]^-\lambda
    & S' \ar[d]^-\lambda \\
  \sF \text{ on }X \ar[r]^-f
    & \dA_S^1 \ar[r]
    & S \text{.}
}\]
One has, on $X_s'=X_{s'}'$,
\[
  \lambda^\ast \eR^i\Psi_\eta(\sF) = \eR^i \Psi_\eta(\sF') = \eR^i\Psi_{\eta'}(\sF')^P \text{:}
\]
this sheaf is constructible, since the induction hypothesis applies to
$X'/S'$, and we use the following lemma applied to $X_s\subset \dA_S^n$ and
to the sheaves of vanishing cycles.




\begin{lemma_}\label{VII:3-5}
Let $\dA^n$ be the affine space of dimension $n$ over a field $k$,
$X\subset \dA^n$, $]sF$ a sheaf of $A$-modules on $X$ and $\bar\eta$ a
generic geometric point of $\dA^1$. Let $X_{\bar\eta,i}$ be the
geometric generic fibre of
$X\subset \dA^n\xrightarrow{\pr_i} \dA^1$ and $\sF_{\bar\eta,i}$ the restriction
of $\sF$ to $X_{\bar\eta,i}$. If the $\sF_{\bar\eta,i}$ are constructible,
there exists $\sF'\subset \sF$, constructible, such that the local sections of
$\sF/\sF'$ are of finite support.
\end{lemma_}

By passage to the limit, for each $i$
\begin{enumerate}[a)]
  \item There exists an \'etale neighbourhood $U$ of $\bar\eta$ and a
    constructible sheaf $\sH$ on $X_{U,i} = X\times_{\dA^1,\pr_i} U$ of restriction
    $\sF_{\bar\eta,i}$ to $X_{\bar\eta,i}$.
  \item For suitable $U$, the isomorphism
    $\sH_{\bar\eta,i} \iso \sF_{\bar\eta,i}$ comes from $u:\sH\to \sF_{U,i}$
    where $\sF_{U,i}$ is the inverse image of $\sF$ on $X_{U,i}$. Let
    $\varphi:X_{U,i} \to X$. The morphism $u$ defines
    $v_i:\varphi_! \sH \to \sF$, of image $\sF_i'$. One has
    $(\sF/\sF_i')_{\bar\eta,i} = 0$ and one takes for $\sF'$ the sum of the
    $\sF_i'$.
\end{enumerate}




\subsection{}\label{VII:3-6}

Let us prove \ref{VII:3-2} in dimension $n$. We may assume $X$ affine (since the
problem is local on $X_s$), then projective over $S$ (take the closure
of $X$ in a projective embedding). One then has a spectral sequence
\cite[I 2.2.3]{sga7} (the Leray spectral sequence for $\bar j$, taking into account the
base-change theorem for the proper morphism $X\to S$).
\begin{equation*}\tag{1}\label{VII:eq:3-6-1}
  E_2^{p q} = \h^p(X_s,\eR^i\Psi_\eta(\sF)) \Rightarrow \h^{p+q}(X_{\bar\eta},\sF) \text{.}
\end{equation*}
Let $\sG^q\subset \eR^q \Psi_\eta(\sF)$ be as in \ref{VII:3-3}, and let
$\sH^q$ be the quotient sheaf. One has $\h^i(X_s,\sH^q)=0$ for $i\ne 0$, and for
$\sH^q$, hence $\eR^q\Psi_\eta(\sF)$, to be constructible, it suffices that
$\h^0(X_s,\sH^q)$ be of finite type.

Let us compute \eqref{VII:eq:3-6-1} modulo finite-type modules (i.e.\ in the
quotient category of that of $A$-modules by the thick subcategory
of $A$-modules of finite type). The long exact cohomology sequence deduced
from $0 \to \sG^q \to \eR^q\Psi_\eta(\sF) \to \sH^q \to 0$ provides, by
\ref{VII:1-10} applied to $\sG^q$, that $E_2^{p q} \sim \h^q(X_s,\sH^q)$.
In particular, $E_2^{p q}\sim 0$ for $p\ne 0$,
$E_2^{0 q}\sim \h^q(X_{\bar\eta},\sF) \sim 0$, and this completes the
proof.

Parallel arguments provide the following result, which improves
\cite[XIII 2.1.12, 2.4.2]{sga7}.




\begin{proposition_}\label{VII:3-7}
The formation of the sheaves of vanishing cycles is compatible with
changes of traits.
\end{proposition_}

Let $g:(S',\eta',s') \to (S,\eta,s)$ be a (surjective) morphism of strictly
local traits, $X/S$, $\sF$ on $X_\eta$, of torsion prime to the
residue characteristic and $(X',\sF')$ their inverse images over
$S'$. Still denoting by $g$ the morphisms parallel to $g$, such as
$X_{s'}' \to X_s$, one must prove that
\begin{equation*}\tag{3.7.1}\label{VII:eq:3-7-1}
\xymatrix{
  g^\ast \eR\Psi_\eta(\sF) \ar[r]^-\sim
    & \eR\Psi_{\eta'}(\sF') \text{.}
}
\end{equation*}
This morphism is defined by the commutative diagram
\[\xymatrix{
  X_{s'}' \ar@{^{(}->}[r]^-{i'} \ar[d]^-g
    & X' \ar[d]^-g
    & X_{\eta'}' \ar[l]_-{\bar j'} \ar[d]^-g \\
  X_s \ar@{^{(}->}[r]^-i
    & X
    & \ar[l]_-{\bar j} X_{\bar\eta}
}\]

A passage to the limit reduces to assuming $X$ of finite type over $S$, and we
proceed by induction on $\dim{X_\eta}$. The question being local, we
may assume $X$ affine, then projective. Assume \eqref{VII:eq:3-7-1} for
$\dim{X_\eta}<n$. For $\dim{X_\eta}=n$, it then follows from the two
following assertions, where $\Delta$ is the mapping cylinder of
\eqref{VII:eq:3-7-1}.

\begin{enumerate}[(A)]
  \item The support of the local sections of the cohomology sheaves of
    $\Delta$ is finite.
    \phantomsection\hypertarget{VII:3-A}{}
  \item $\eR\Gamma(X_s,\Delta)=0$, hence $\Delta=0$, by
    (\hyperlink{VII:3-A}{A}).
    \phantomsection\hypertarget{VII:3-B}{}
\end{enumerate}




\subsection{}\label{VII:3-8}

The proof of (\hyperlink{VII:3-A}{A}) is parallel to that of
\ref{VII:3-3}. We localize on $X$ and consider maps
$X\to \dA_S^1$; denoting by $S(x)$ the strict henselization of $\dA_S^1$ at the
generic point of the special fibre, we apply the induction hypothesis
to $X_1/S(x)$ ($X_1=X\times_{\dA^1} S(x)$) and to $S'(x)$. Taking
invariants under a pro-$p$-group, we find that $\Delta$ is zero on the
geometric generic fibre of $X_{s'} \to \dA_s^1$. This being true,
locally on $X$ and for every projection, we have (\hyperlink{VII:3-A}{A}).




\subsection{}\label{VII:3-9}

For (\hyperlink{VII:3-B}{B}), properness of $X$ over $S$ ensures that
\[
  \eR\Gamma(X_{s'}',\eR\Psi_{\eta'}(\sF')) = \eR\Gamma(X_{\bar\eta}',\sF')
\]
and
\[
  \eR\Gamma(X_s,\eR\Psi_\eta(\sF)) = \eR\Gamma(X_{\bar\eta},\sF) \text{.}
\]
One has
$\eR\Gamma(X_s,\eR\Psi_\eta(\sF))\iso \eR\Gamma(X_{s'}',g^\ast\eR\Psi_\eta(\sF))$
(invariance under change of separably closed base field);
$\eR\Gamma(X_{s'}', \Delta)$ is therefore the mapping cylinder of
$\eR\Gamma(X_{\bar\eta},\sF) \to \eR\Gamma(X_{\bar\eta}',\sF')$, an
isomorphism by the same invariance theorem.




\subsection{}\label{VII:3-10}

Let us prove \ref{VII:1-1}. We reduce to assuming $S$ connected, hence
integral. If $\dim S=0$, ($S$ spectrum of a field); we apply \ref{VII:1-9},
\ref{VII:1-10}. If $S$ is of dimension $1$, theorem \ref{VII:1-9}
ensures that the $\eR^i f_\ast \sF$ are constructible above the complement
of a finite set $T$ of points of $S$. It remains to see that their restriction
to the $Y_t$ ($t\in T$) are constructible. It suffices to see this after
localization at $t$: we may assume that $S$ is a strictly local trait.
The essential case is the following.




\begin{lemma_}\label{VII:3-11}
Let $X$ of finite type over $S$, $j:X_\eta\hookrightarrow X$,
$i:X_s\hookrightarrow X$ and $\sF$ constructible on $X_\eta$. Then
$i^\ast\eR j_\ast \sF$ is constructible.
\end{lemma_}

Let $I=\gal(\bar\eta/\eta)$ be the inertia group. We know that it is an
extension of $\widehat\dZ_{p'}(1)$ by a pro-$p$-group $P$ ($p$ the
characteristic exponent of the residue field). The Hochschild-Serre spectral sequence
gives
\[
  E_2^{p q} = \sH^p(I,\eR^q\Psi_\eta(\sF))\Rightarrow i^\ast \eR^{p+q}j_\ast\sF \text{.}
\]
Let $R_t^q = \eR^q\Psi_\eta(\sF)^P$ ($t$ for tame). If $\sigma$ is a
generator of $\widehat\dZ_{p'}(1)$, one has
$E_2^{0 q} = \ker(\sigma-1,R_t^q)$,
$E_2^{1 q} = \operatorname{coker}(\sigma-1,R_t^q)$ and $E_2^{p q} = 0$ for
$p\ne 0,1$. These sheaves are constructible, since the sheaves of vanishing
cycles are, and \ref{VII:3-11} follows.




\subsection{}\label{VII:3-12}

Let us treat the general case, $S$ still being a strictly local
henselian trait. If $X=X_\eta$, $f$ is the composite $X\to Y_\eta \to Y$ and we
apply \ref{VII:1-9} and \ref{VII:3-11}. In the general case, let
$j:X_\eta\hookrightarrow X$ and let $\Delta$ be the mapping cylinder of
$\sF \to \eR j_\ast j^\ast \sF$. Its cohomology sheaves are supported
on $X_s$; by \ref{VII:1-9} applied to $X_s \to Y_{s'}$,
$\eR f_\ast \Delta$ is therefore constructible. The complex
$\eR f_\ast \eR j_\ast j^\ast \sF = \eR(f j)_\ast j^\ast \sF$ is
also, and we conclude by the triangle
\[\xymatrix{
  \ar[r]
    & \eR f_\ast \sF \ar[r]
    & \eR f_\ast \eR j_\ast j^\ast \sF \ar[r]
    & \eR f_\ast \Delta \ar[r]
    & \text{.}
}\]










\section{Local biduality}\label{VII:4}




\subsection{}\label{VII:4-1}

Let $S$ be regular of dimension $0$ or $1$ and put $K_S=A$ (constant
sheaf with value $A$). For $K\in\ob\eD_\text{ctf}^b(S,A)$, put
$D K=\rHom(K,K_S)\in\ob\eD^+(S,A^\circ)$ (a complex of sheaves of
right $A$-modules). One still has $D K\in\ob\eD_\text{ctf}^b(S,A^0)$, and
an explicit computation, possible since $S$ is of dimension $1$, shows
$K\iso D D K$ (for $A=\dZ/n$, \hyperref[V]{Duality} \ref{V:1-4}).




\subsection{}\label{VII:4-2}

We shall assume that $A$ is commutative, an hypothesis probably unnecessary. For
$a:X\to S$ of finite type over $S$, put $K_X=\eR a^! K_S$. For
$K\in\ob\eD_\text{ctf}^b(S,A)$, put
$D K=\rHom(K,K_X)\in\ob\eD_\text{ctf}^b(X,A)$. Our main result is
the following




\begin{theorem_}\label{VII:4-3}
One has $K\iso D D K$ ($K\in \ob\eD_\textnormal{ctf}^b(X,A)$).
\end{theorem_}




\begin{lemma_}\label{VII:4-4}
If $a$ is proper, one has $\eR a_\ast K \iso \eR a_\ast D D K$.
\end{lemma_}

Poincar\'e duality
\[
  \eR a_! \rHom(K,\eR a^! K_S) = \rHom(\eR a_! K,K_S)
\]
provides us with an isomorphism $\eR a_\ast D = D \eR a_! = D \eR a_\ast$. The
lemma follows from commutativity of the diagram (\hyperref[V]{Duality},
\ref{V:1-2})
\[\xymatrix{
  \eR a_\ast K \ar[r] \ar[d]^-\wr
    & \eR a_\ast D D K \ar[d]^-\wr \\
  D D \eR a_\ast K \ar[r]^-\sim
    & D \eR a_\ast D K \text{.}
}\]




\subsection{}\label{VII:4-5}

We see by localization that it suffices to treat the cases where $S$ is the spectrum
of a field, or a trait. For $S$ the spectrum of a field, we proceed by
induction on $\dim X$. The problem being local, we may assume $X$
proper. Moreover, the induction hypothesis ensures, by the usual
arguments, that the mapping cylinder $\Delta$ of $K\to D D K$ has
skyscraper cohomology sheaves. The lemma then ensures that $\eR a_\ast \Delta=0$,
hence that $\Delta=0$.




\subsection{}\label{VII:4-6}

For $S$ a trait, we already know that \ref{VII:4-3} holds on the generic
fibre (\ref{VII:4-5}). We may also assume $X$ proper over $S$ and we
proceed by induction on $\dim{X_s}$, $X$ proper over $S$. The mapping
cylinder $\delta$ again has skyscraper cohomology sheaves, on
$X_s$, and we conclude as above.




\subsection{}\label{VII:4-7}

Assume that $A=\dZ/m$. One then has a variant of the preceding
theory, taking $K$ arbitrary in $\eD_c^b(X,A)$. The complex
$\eR a^! \dZ/m$ is of finite injective dimension \cite[XVIII 3.1.7]{sga4}. This
ensures that $D K$ is still in $\eD_c^b(X,A)$. To prove local biduality
over $S$, one uses that $\dZ/m$ is the dualizing $\dZ/m$-module.
Next, one proceeds as above.










\section{Appendix, by L.\ Illusie}\label{VII:5}

In this appendix, which resumes certain parts of the late \cite[II]{sga5}, and
a letter from P.\ Deligne to the author, we prove \ref{VII:2-14} as well as
various generalizations and variants of the specialization theorems
\cite[XVI 2.1]{sga4} and (\hyperref[I]{\'Etale cohomology} \ref{I:5-1-7}).

The schemes and morphisms considered will be assumed quasi-compact and
quasi-separated. If $X$ is a scheme and $x$ a geometric point of
$X$, we shall denote by $X(x)$ the strict localization of $X$ at $x$.

We fix a base scheme $S$ and a sheaf of rings $A$ over $S$ (not
necessarily commutative nor noetherian). If $X$ is an $S$-scheme, we
shall write $\eD(X)$ for $\eD(X,A_X)$.




\subsection{Cohomological properness}\label{VII:5-1}




\begin{proposition}\label{VII:5-1-1}
Let $f:X\to Y$ be an $S$-morphism and $E\in\ob \eD^+(X)$. The following
conditions are equivalent:
\begin{enumerate}[\indent (i)]
  \item The formation of $\eR f_\ast E$ commutes with every finite base
    change $S'\to S$.
  \item The formation of $\eR f_\ast E$ commutes with every quasi-finite base
    change (or projective limit of quasi-finite morphisms) $S'\to S$.
  \item For every specialization arrow $i:t\to S(s)$
    (\hyperref[I]{\'Etale cohomology}, \ref{I:5-1-2}), if we denote by
    $\widetilde S$ the normalization in $k(\bar t)$ of the integral scheme
    closure of $i(t)$ in $S(s)$ (by
    \cite[IV, 6.15.5, 18.6.11, 18.8.16]{ega}) $\bar S$ is thus local
    integral, with closed point radicial over $s$), $\bar f:\bar X\to \bar Y$
    the morphism deduced from $f$ by the base change $\bar S\to S$,
    $f_s:X_s\to Y_s$ the fibre of $f$ at $s$, then the base-change arrow
    \[
      (\eR \bar f_\ast (E|\bar X))|Y_s \to \eR {f_s}_\ast(E|X_s)
    \]
    is an isomorphism.
\end{enumerate}
\end{proposition}

The equivalence of (i) and (ii) follows from the main theorem of Zariski,
and it is clear that (ii) implies (iii). Let us prove that (iii) implies (ii). For
every $S'\to S$, projective limit of quasi-finite morphisms, and every geometric
point $s$ of $S'$, denote by $C(S',s)$ the cone of the base-change
arrow
\[
  (\eR f_\ast'(E|X'))|Y_s \to \eR{f_s}_\ast (E|X_s) \text{,}
\]
where $f':X'\to Y'$ is the morphism deduced from $f$ by the base change
$S'\to S$ and $f_s:X_s\to Y_s$ the fibre of $f$ at $s$. We must prove that,
for every $(S',s)$ as above, $C(S',s)$ is acyclic. We shall show, by
induction on $N$, that $\h^i C(S',s) = 0$ for $i\leqslant N$. By the Main
Theorem, passage to the limit, and elimination of radicial morphisms, we may
assume that $S$ and $S'$ are strictly local, $S'$ integral, with closed
point $s$, $S'\to S$ finite. Let $t$ be a generic geometric point
of $S'$, and $\bar S$ the normalization of $S'$ in $k(t)$. Let $\bar S_\bullet$
be the coskeleton simplicial scheme of $S'\to S'$
($\bar S_0=\bar S$,\ldots $\bar S_n=(\bar S/S')^{n+1}$,\ldots),
$X_\bullet=X\times_S \bar S_\bullet$. $g_\bullet:\bar X_\bullet \to X'$ the
canonical projection. Put $G|X'=G'$. The bicomplex
${g_\bullet}_\ast g_\bullet^\ast G'$, with columns the
${g_n}_\ast g_n^\ast G'$ is a resolution of $G'$ \cite[VIII 8]{sga4}. One
can therefore represent $C(S',s)$ by a bicomplex whose columns
are identified with the $C(\bar S_n,s)$. By (iii), $C(\bar S_0,s)=C(\bar S,s)$
is acyclic. On the other hand, for $n>0$, one has $\h^i(\bar S_n,s)=0$ for
$i\leqslant N$ by the induction hypothesis. We conclude that
$\h^i C(S',s)=0$ for $i\leqslant N+1$, which completes the proof.




\subsubsection{Remark}\label{VII:5-1-2}

Let $(f,E)$ satisfy condition (i) of \ref{VII:5-1-1}.
\begin{enumerate}[a)]
  \item If $E\in \ob \eD^b(X)$, and $f_\ast$ is of finite cohomological
    dimension, then, for every $F\in \ob\eD^-(S,A^\circ)$, the canonical arrow
    \begin{equation*}\tag{$*$}\label{VII:eq:5-1-1-1}
      \eR f_\ast(E)\lotimes_A q^\ast F \to \eR f_\ast(E\lotimes_A p^\ast F) \text{,}
    \end{equation*}
    where $p:X\to S$, $q:Y\to S$ are the projections, is an isomorphism.
    Same conclusion with $E\in\ob\eD^-(X)$.
\end{enumerate}

By d\'evissage, it indeed suffices to verify the assertion for $F$ of the
form $i_! A_{S'}$, with $i:S'\to S$ quasi-finite. By the Main Theorem, we reduce
to assuming $i$ finite, the conclusion is then equivalent to the
commutation of $\eR f_\ast(E)$ with base change by $i$.
\begin{enumerate}[a)]
\setcounter{enumi}{1}
  \item Assume $A$ noetherian, constructible. Then, for every
    $F\in\ob\eD^b(S,A^\circ)$, constructible and of finite tor-dimension, the
    arrow \eqref{VII:eq:5-1-1-1} above is an isomorphism.
\end{enumerate}

We reduce indeed, by d\'evissage, to the case where $F$ is of the form
$i_! M$, with $i:S'\to S$ locally closed, $A$ locally constant on $S'$,
$M$ constructible, of finite tor-dimension, and with locally constant
cohomology on $S'$. Changing base from $S$ to $S'$, we are thus reduced to the
case where $F$ is perfect \cite[I]{sga6}, hence finally, by localization and
d\'evissage, to the case where $F=A$, for which the conclusion is trivial.




\subsubsection{}\label{VII:5-1-3}

Let $f:X\to Y$ be an $S$-morphism, and $E\in\ob\eD^+(X)$. We shall say that
$(f,E)$ is \emph{cohomologically proper} rel.\ $S$ if the formation of
$\eR f_\ast E$ commutes with every base change $S'\to S$. It amounts to the
same to say that after every base change $S'\to S$, the formation of
$\eR f_\ast' E'$ (where $f'=f\times_S S'$, $E'=$inverse image of $E$ on
$X\times_S S'$) commutes with every finite base change $S'' \to S'$ (cf.\
\cite[XII 6.1]{sga4}). Here are some examples of cohomological properness
(for a study of the analogous notion for sheaves of sets or of
non-commutative groups, the reader is referred to \cite[XIII]{sga1}).
\begin{enumerate} % originally 1.3.1, 1.3.2, ...
  \item Assume $A$ torsion and $f$ proper. Then, for every
    $E\in\ob\eD^+(X)$, $(f,E)$ is cohomologically proper rel.\ $S$
    (proper base-change theorem, \cite[XII 5.1]{sga4}).
  \item Assume $S$ finite, $X$ and $Y$ of finite type over $S$, and $A$ annihilated
    by an integer invertible on $S$. Then, for every $E\in\ob\eD^+(X)$,
    $(f,E)$ is cohomologically proper rel.\ $S$
    (\ref{VII:1-9}).
  \item Assume $A$ constant, noetherian, annihilated by an integer $n$
    invertible on $S$, $f:X=Y\setminus D\to Y$ the inclusion of the complement
    of a divisor $D$ on $Y$, with normal crossings rel.\ $S$
    \cite[XIII 2.1]{sga1}. Let $E$ be a sheaf of $A$-modules on $X$
    satisfying one of the following conditions:
    \begin{enumerate}[(i)]
      \item $E$ is locally constant constructible and tamely
        ramified along $D$;
      \item $E$ is locally for the \'etale topology on $Y$ and on $S$
        the inverse image of a sheaf of $A$-modules on $S$.
    \end{enumerate}
    Then $(f,E)$ is cohomologically proper rel.\ $S$, and $\eR f_\ast E$ is
    constructible (compare with \cite[XIII 2.4]{sga1}).
\end{enumerate}

Let us sketch the proof in case (i). The question is local in a
neighbourhood of a point $y$ of $Y$. We may assume $D$ a sum of smooth
divisors, $D=\sum_{i=1}^m D_i$. By virtue of the relative Abhyankar lemma
\cite[XIII 5.5]{sga1}, we may, in a neighbourhood of $y$, trivialize $E$ by a
finite covering $\widetilde Y\to Y$ of the form
$\widetilde Y=Y[T_1,\dots,T_r]/(T_1^{n_1}-t_1,\dots,T_r^{n_r}-t_r)$ where the
$t_i$ are local equations of the smooth divisors passing through $y$, and the
$n_i$ integers prime to the char.\ of $k(y)$. The inverse image
$\widetilde D$ of $D$ in $Y$ is still with normal crossings rel.\ $S$,
and if $g:\widetilde X=\widetilde Y\setminus \widetilde D\to X$ is the
projection, $g^\ast E$ is constant. Since $E$ injects into $g_\ast g^\ast E$
and the quotient is tamely ramified, an easy d\'evissage reduces
to the case where $E$ is constant. Denote by $p:X\to S$, $q:Y\to S$ the
projections, and, for $1\leqslant i\leqslant m$, $f_i:Y\setminus D_i \to D$
the canonical inclusion. From the relative purity theorem \cite[XVI 3.7]{sga4}
we easily deduce, by induction on $m$, that, for every locally constant
$A$-module $M$ on $S$, the canonical arrow
\begin{equation*}\tag{5.1.3.1}\label{VII:eq:5-1-3-3-1} % originally 1.3.3.1
  \eR f_\ast (\dZ/n)\lotimes q^\ast M\to \eR f_\ast (p^\ast M)
\end{equation*}
is an isomorphism, and that on the other hand one has:
\begin{align*}\tag{5.1.3.2}\label{VII:eq:5-1-3-3-2}
  {f_i}_\ast(\dZ/n) &= \dZ/n \text{,} \\
  \eR^1 {f_i}_\ast (\dZ/n) &= (\dZ/n)(-1)_{D_i} \text{,}
\end{align*}
a canonical isomorphism being given by the fundamental class $\cl(D_i)$
(\hyperref[IV]{Cycle} \ref{IV:2-1-4}),
\begin{align*}
  \eR^q {f_i}_\ast(\dZ/n) &= 0\text{ for $q>1$,} \\
  \eR f_\ast (\dZ/n) &= \bigotimes_i^\eL \eR {f_i}_\ast (\dZ/n) \text{ , i.e.:} \\
  \eR^1 f_\ast(\dZ/n) &= \bigoplus_{i=1}^m (\dZ/n)(-1)_{D_i} \text{,} \\
  \textstyle\bigwedge^\bullet \eR^1 f_\ast(\dZ/n) &\iso \eR^\bullet f_\ast (\dZ/n) \text{.}
\end{align*}
The conclusion of 3, in case (i), therefore follows from
\eqref{VII:eq:5-1-3-3-1} and \eqref{VII:eq:5-1-3-3-2}. In case (ii), we reduce,
by passage to the limit, to the case where $E=p^\ast M$, with $M$
constructible, then, by d\'evissage and base change (using (i)), to the
case $M$ constant, already treated.

It follows from the preceding proof that \eqref{VII:eq:5-1-3-3-1}
is an isomorphism for every $M\in\ob\eD^+(S)$.

We shall compare the formulas \eqref{VII:eq:5-1-3-3-2} with the analogous formulas in
integral cohomology in the case $S=\spec(\dC)$ (see for example
\cite[3.1]{de71-2}), which follow from the fact that, locally for the
classical topology, the complement of $D$ has the homotopy type of a torus.

By passage to the limit, we deduce from \eqref{VII:eq:5-1-3-3-2} similar formulas
in $\ell$-adic cohomology, i.e.\ with $\dZ/n$ replaced by
$\dZ_\ell$ ($\ell$ a prime number invertible on $S$). Suppose
$q:Y\to S$ proper and smooth. It then follows from the Weil conjectures
\cite{de74,de80} that the Leray spectral sequence
\[
  E_2^{i j} = \eR^i q_\ast \eR^j f_\ast \dZ_\ell \Rightarrow \eR^\bullet p_\ast \dZ_\ell
\]
degenerates at $E_3$ modulo torsion \cite[\S 6]{de71-1}. The same
phenomenon occurs in integral cohomology in the case
$S=\spec(\dC)$, see (loc.\ cit.) and \cite[3.2.13]{de71-2}.




\subsection{Specialization and cospecialization}\label{VII:5-2}

Let $f:X\to S$ and $E\in\ob\eD^+(X)$.




\subsubsection{}\label{VII:5-2-1}

Assume that the formation of $\eR f_\ast E$ commutes with every finite base
change. For every specialization $u:t\to S(s)$ of geometric points
of $S$, we then define an arrow, called the \emph{specialization
arrow},
\begin{equation*}\tag{5.2.1.1}\label{VII:eq:5-2-1-1}
  \spe(u):\eR\Gamma(X_s,E|X_s) \to \eR\Gamma(X_t,E|X_t) \text{,}
\end{equation*}
in the following manner. Let, as in \ref{VII:5-1-1}(iii), $\bar S$ be the
normalization in $k(t)$ of the closure of $u(t)$ in $S(s)$, and
$\bar f:\bar X\to \bar S$ the morphism deduced from $f$ by the base
change $\bar S\to S$. By the trivial part (i)$\Rightarrow$(iii) of
\ref{VII:5-1-1}, the base-change arrow
\begin{equation*}\tag{5.2.1.2}\label{VII:eq:5-2-1-2}
  \eR\Gamma(\bar X,E|\bar X)=\eR\bar f_\ast (E|\bar X)_s \to \eR\Gamma(X_s,E|X_s)
\end{equation*}
is an isomorphism. We define \eqref{VII:eq:5-2-1-1} as the composite of
the inverse of \eqref{VII:eq:5-2-1-2} and the restriction arrow
\begin{equation*}\tag{5.2.1.3}\label{VII:eq:5-2-1-3}
  \eR\Gamma(\bar X,E|\bar X) \to \eR\Gamma(X_t,E|X_t) \text{.}
\end{equation*}




\subsubsection{}\label{VII:5-2-2}

Assume that $f$ is locally acyclic rel.\ $E$ (\ref{VII:2-12}).
For every specialization $u:t\to S(s)$ of geometric points of $S$, we
then define an arrow, called the \emph{cospecialization arrow},
\begin{equation*}\tag{5.2.2.1}\label{VII:eq:5-2-2-1}
  \cosp(u):\eR\Gamma(X_t,E|X_t) \to \eR\Gamma(X_s,E|X_s) \text{,}
\end{equation*}
defined in the following manner. Let $g:X_t\to X$, $g':\bar X\to X$ be the
canonical arrows. It follows from the local acyclicity hypothesis that
the canonical arrow
\begin{equation*}\tag{5.2.2.2}\label{VII:eq:5-2-2-2}
  \eR g_\ast'(E|\bar X)\to \eR g_\ast(E|X_t)
\end{equation*}
is an isomorphism. Indeed, we may assume $S$ strictly local,
integral, with closed point $s$, $t$ being a generic
geometric point. Let us compute the fibre of \eqref{VII:eq:5-2-2-2} at a
geometric point $x$ of $X$, of image $y$ in $S$. After replacing $S$ by
the strict localization of $S$ at $y$, we may assume that $y=s$. Since
$\bar S\to S$ is a projective limit of finite morphisms, and radicial at $s$, we
find that $\eR g_\ast'(E|\bar X)_x = E_x$ and that the fibre of
\eqref{VII:eq:5-2-2-2} at $x$ identifies with the restriction
$E_x=\eR\Gamma(X(x),E)\to \eR\Gamma(X(x)_t,E|X(x)_t)$, which is an
isomorphism by hypothesis. Applying $\eR\Gamma(X,-)$ to
\eqref{VII:eq:5-2-2-2}, we deduce that the restriction arrow
\eqref{VII:eq:5-2-1-3} is an isomorphism. We define \eqref{VII:eq:5-2-2-1}
as the composite of the inverse of \eqref{VII:eq:5-2-1-3} and the
base-change arrow \eqref{VII:eq:5-2-1-2}.




\subsubsection{}\label{VII:5-2-3}

If $s'' \xrightarrow{u'} s' \xrightarrow u s$ are specialization
arrows of geometric points of $S$, one checks that one has
\begin{align*}
  \spe(u')\spe(u'') &= \spe(u u') \text{,} \\
  \cosp(u) \cosp(u') &= \cosp(u u') \text{,}
\end{align*}
whenever both sides are defined.




\subsubsection{}\label{VII:5-2-4}

Assume that the formation of $\eR f_\ast E$ commutes with every finite base
change, and that $f$ is locally acyclic rel.\ $E$. It follows from the
definitions that, for every specialization $u:t\to S(s)$ of geometric
points of $S$; the arrows $\spe(u)$ and $\cosp(u)$ are
isomorphisms, inverse to each other.

The preceding hypotheses are satisfied in particular when $A$ is
noetherian, annihilated by an integer $n$ invertible on $S$, $f$ proper and
smooth, and $E$ with locally constant cohomology (cohomological properness
of proper morphisms (\ref{VII:5-1-3}), and local acyclicity of smooth morphisms
\cite[XV]{sga4}). We recover the ``specialization theorem''
\cite[XVI 2.1]{sga4}.

On the other hand, \ref{VII:5-1-1} has the following consequence:




\begin{proposition}\label{VII:5-2-5}
Assume that $f$ is locally acyclic relative to $E$, and that, for
every specialization $u:t\to S(s)$, the cospecialization arrow
$\cosp(u)$ is an isomorphism. Then the formation of $\eR f_\ast E$ commutes
with every finite base change.
\end{proposition}

Taking into account \ref{VII:5-2-4}, \ref{VII:5-2-5} generalizes
(\hyperref[I]{\'Etale cohomology} \ref{I:5-1-7}).




\begin{corollary}\label{VII:5-2-6}
Assume that $f$ is locally acyclic rel.\ $E$. Let $x$ be a geometric
point of $X$, of image $s$ in $S$, denote by $f_{(x)}:X(x)\to S(s)$ the
morphism induced by $f$. Then the formation of $\eR {f_{(x)}}_\ast(E|X(x))$
commutes with every finite base change $S'\to S(s)$.
\end{corollary}

It suffices, by \ref{VII:5-2-5}, to check that the
cospecialization arrows for $(f_{(x)},E|X(x))$ are isomorphisms. By
transitivity of the cospecialization arrows (\ref{VII:5-2-3}), it suffices
to show that, if $u:t\to S(s)$ is a specialization, then
$\cosp(u):\eR\Gamma(X(x)_t,E|X(x)_t)\to \eR\Gamma(X(x)_s,E|X(x)_s)$ is an
isomorphism. But, since $s$ is closed in $S(s)$, $X(x)_s$ is strictly
local with closed point $x$, $\eR\Gamma(X(x)_s,E|X(x)_s)=E_x$, and
$\cosp(u)$ is an isomorphism, inverse of the restriction isomorphism
$E_x\to \eR\Gamma(X(x)_t,E|X(x)_t)$.




\begin{corollary}\label{VII:5-2-7}
Let $f:X\to Y$ be an $S$-morphism, $p:X\to S$, $q:Y\to S$ the projections, and
$E\in\ob\eD^+(X)$. Assume $A$ locally constant, $f$ locally acyclic
rel.\ $E$, and $q$ locally acyclic at every geometric point $y$ of
$Y$ rel.\ every sheaf locally constant in a neighbourhood of $y$. Then $p$
is locally acyclic rel.\ $E$.
\end{corollary}

Let $x$ be a geometric point of $x$, of images $y$ and $s$ in $Y$ and
$S$, $t\to S(s)$ a specialization, $f_{(x)}:X(x)\to Y(y)$ the morphism
induced by $f$. We must show that the restriction
\begin{equation*}\tag{$*$}\label{VII:eq:5-2-7-1}
  E_x \to \eR\Gamma(X(x)_t,E|X(x)_t)
\end{equation*}
is an isomorphism. We may assume $A$ constant. By \ref{VII:5-2-4} and
\ref{VII:5-2-5}, for every specialization $z\to Y(y)$, the
specialization arrow
\[
  E_x = \eR {f_{(x)}}_\ast (E|X(x))_y \to \eR {f_{(x)}}_\ast (E|X(x)|_z = \eR\Gamma(X(x)_z,E|X(x)_z)
\]
is an isomorphism, in other words $\eR{f_{(x)}}_\ast(E|X(x))$ is constant of
value $E_x$. The acyclicity hypothesis on $q$ therefore implies that the
restriction
\[
  E_x \to \eR\Gamma(Y(y)_y,(E_x)|Y(y)_t) = \eR\Gamma\left(Y(y)_t,\eR {f_{(x)}}_\ast (E|X(x))|Y(y)_t\right)
\]
is an isomorphism. But this arrow identifies with
\eqref{VII:eq:5-2-7-1}, whence the conclusion.

The hypothesis of \ref{VII:5-2-7} on $q$ is satisfied in particular when $q$
is smooth, and $A$ is annihilated by an integer $n$ invertible on $S$
\cite[XV]{sga4}. On the other hand, if the hypotheses of \ref{VII:5-2-7} are
satisfied after every base change $S'\to S$, we then conclude that
$p$ is universally locally acyclic rel.\ $E$. From these remarks
follows in particular the statement \ref{VII:2-14}.




\begin{corollary}\label{VII:5-2-8}
Let $f:X\to S$ and $E\in \ob\eD^+(X)$. Assume that $f$ is universally
locally acyclic rel.\ $E$, that the formation of $\eR f_\ast E$ commutes
with every finite base change, and that for every morphism $s'\to s$, where
$s$ is a geometric point of $S$ and $s'$ the spectrum of a
separably closed field, the base-change arrow
$\eR\Gamma(X_s,E|X_s)\to \eR\Gamma(X_{s'},E|X_{s'})$ is an isomorphism.
Then $(f,E)$ is cohomologically proper (\ref{VII:5-1-3}).
\end{corollary}

Let $f':X'\to S'$ be the morphism deduced from $f$ by a base change
$S'\to S$, we must see that the formation of $\eR f_\ast'(E|X')$ commutes
with every finite base change. By \ref{VII:5-2-5}, it suffices to
show that, for every specialization $u':t'\to S'(s')$, the
cospecialization arrow $\cosp(u')$:
$\eR\Gamma(X_{t'}',E|X_{t'}') \to \eR\Gamma(X_{s'}',E|X_{s'}')$ is an
isomorphism. Let $s$ (resp.\ $t$) be the geometric point image of $s$
(resp.\ $t$) in $S$, $u:t\to S(s)$ the specialization defined by $u'$.
One easily sees that one has a commutative square
\[\xymatrix{
  \eR\Gamma(X_t,E|X_t) \ar[r]^-{\cosp(u)} \ar[d]
    & \eR\Gamma(X_s,E|X_s) \ar[d] \\
  \eR\Gamma(X_{t'}',E|X_{t'}') \ar[r]^-{\cosp(u')}
    & \eR\Gamma(X_{s'}',E|X_{s'}')
}\]
where the vertical arrows are the base-change arrows. By
hypothesis, these are isomorphisms. On the other hand (\ref{VII:5-2-4})
$\cosp(u)$ is an isomorphism, so $\cosp(u')$ is an isomorphism, which
completes the proof.

Here is an example of application of \ref{VII:5-2-8}. Assume $A$ annihilated by
an integer $n$ invertible on $S$, $f$ locally of finite type and
universally locally acyclic. Let $x$ be a geometric point of
$X$, of image $s$ in $S$, $f_{(x)}:X(x)\to S(s)$ the morphism induced by $f$.
It follows from \ref{VII:5-2-6} and \ref{VII:1-9} (by passage to the limit)
that the hypotheses of \ref{VII:5-2-8} are satisfied by the pair
$(f_{(x)},E|X(x))$, which is therefore cohomologically proper.




\subsubsection{}\label{VII:5-2-9}

Assume $A$ constant. Let $f:X\to S$, and $E\in\ob\eD^+(X)$ a complex of
finite tor-dimension. We shall say that $f$ is \emph{strongly locally
acyclic} rel.\ $E$ if, for every geometric point $x$ of $X$,
of image $s$ in $S$, every specialization $t\to S(s)$, and every
$A^\circ$-module $M$, the restriction arrow
\[
  E_x\lotimes_A M \to \eR\Gamma(X(x)_t,(E|X(x)_t)\lotimes_A M)
\]
is an isomorphism. I do not know whether ``locally acyclic'' implies
``strongly locally acyclic.''




\begin{proposition}\label{VII:5-2-10}
Under the hypotheses of \ref{VII:5-2-9}, assume $A$ noetherian of
torsion, $S$ noetherian, and $f$ strongly locally acyclic rel.\ $E$.
Then, for every diagram with cartesian squares
\[\xymatrix{
  X \ar[d]^-f
    & \ar[l]_-j X' \ar[d]^-{f'}
    & \ar[l]_-{j'} X'' \ar[d]^-{f''} \\
  S
    & \ar[l]_-i S'
    & \ar[l]_-{i'} S'' \text{,}
}\]
with $i$ quasi-finite and $i'$ an open immersion, and every
$F\in\ob\eD^+(S'',A^\circ)$, the canonical arrow
\begin{equation*}\tag{$*$}\label{VII:eq:5-2-10-1}
  j^\ast E \lotimes {f'}^\ast \eR i_\ast' F \to \eR j_\ast' ((j j')^\ast E\lotimes {f''}^\ast F)
\end{equation*}
is an isomorphism.
\end{proposition}

We shall compare this statement with \cite[XV 1.17]{sga4}.

Let us prove \ref{VII:5-2-10}. We easily reduce to $S=S'$,
strictly local, and $F$ constructible, concentrated in degree $0$. We
then resolve $F$ on the right by finite sums of sheaves of the form
$u_\ast M$, where $u:t\to S$ is a geometric point and $M$ an
$A^\circ$-module. The hypothesis implies that \eqref{VII:eq:5-2-10-1} is an
isomorphism for $F=u_\ast M$, whence the result.
